\subsection{TimeMemEval Benchmark Construction}
  \label{app:timememeval-construction}

  \begin{tcolorbox}[
    enhanced jigsaw,
    breakable,
    colback=white,
    colframe=black,
    boxrule=0pt,
    borderline north={0.7pt}{0pt}{black},
    borderline south={0.7pt}{0pt}{black},
    sharp corners,
    left=0pt,
    right=0pt,
    top=2pt,
    bottom=2pt,
    before skip=0.6em,
    after skip=0.8em,
  ]
  \scriptsize
  \noindent\textbf{Algorithm 2} \textsc{TimeMemEval} Benchmark Construction Procedure

  \vspace{0.15em}
  \hrule
  \vspace{0.45em}

  \noindent\textbf{Require:}\\
  \hspace*{1.5em}Persona/theme pools, event-family library, bridge scopes
  $\mathcal{B}$, offset sampler $\Omega$, terminal types $\mathcal{Y}$\\
  \hspace*{1.5em}Distractor modes $\mathcal{M}$, surface temporal styles $\mathcal{R}$,
  event/session generators, uniqueness judge $\mathcal{J}$, temporal validator
  $\mathcal{V}$, retry budgets

  \vspace{0.25em}
  \noindent\textbf{Ensure:}\\
  \hspace*{1.5em}\textsc{TimeMemEval} records with \texttt{record.json},
  \texttt{sessions.json}, generation traces, split metadata, and validation logs

  \vspace{0.25em}
  \begin{tabular}{@{}r@{\quad}p{0.84\linewidth}@{}}
  1: & \textbf{function} BuildTimeMemEval$(\mathcal{B},\Omega,\mathcal{Y})$ \\
  2: & \hspace*{1em}\textbf{for each} requested record \textbf{do} \\
  3: & \hspace*{2em}Sample persona/theme context and generate a source event $e_s$ with a valid temporal anchor \\
  4: & \hspace*{2em}Sample bridge scope and offset $(\beta,\delta)$, where $\beta\in\mathcal{B}$ and $\delta\in\Omega$ \\
  5: & \hspace*{2em}$W^\star \leftarrow$ ComputeTargetWindow$(e_s.\mathrm{event\_time},\beta,\delta)$ \\
  6: & \hspace*{2em}Sample terminal type $y\in\mathcal{Y}$ and a compatible event family; generate destination event $e_d$ inside $W^\star$ \\
  7: & \hspace*{2em}Generate probe question $q$ and answer $a$ that require linking $e_s$ to $e_d$ through $W^\star$ \\
  8: & \hspace*{2em}Generate hard distractors $\mathcal{D}$ using modes in $\mathcal{M}$, including same-activity, answer-alias, near-window, and source-anchor-alias distractors outside $W^\star$ \\
  9: & \hspace*{2em}Generate filler events $\mathcal{U}$ that add dialogue noise without leaking the answer or temporal bridge \\
  10: & \hspace*{2em}\textbf{if} $\mathcal{J}(e_s,e_d,\mathcal{D},\mathcal{U},q,a)$ rejects uniqueness or answerability \textbf{then} retry the failed stage \\
  11: & \hspace*{2em}\textbf{for each} event $e\in\{e_s,e_d\}\cup\mathcal{D}\cup\mathcal{U}$ \textbf{do} \\
  12: & \hspace*{3em}Sample a surface temporal style $r\in\mathcal{R}$ and generate a multi-turn session $s_e$ with relative, indirect, or calendar-based temporal phrasing \\
  13: & \hspace*{3em}\textbf{if} $\mathcal{V}(s_e,e,\beta)$ rejects temporal consistency \textbf{then} regenerate $s_e$ or its dependent event \\
  14: & \hspace*{2em}\textbf{end for} \\
  15: & \hspace*{2em}Write record, sessions, traces, split metadata, and validation logs \\
  16: & \hspace*{1em}\textbf{end for} \\
  17: & \textbf{end function}
  \end{tabular}
  \end{tcolorbox}

  \textsc{TimeMemEval} is designed as a controlled benchmark for temporal-bridge
  retrieval in long-term memory, rather than covering all forms of
  temporal reasoning. Each instance requires a model to identify a source event,
  derive a target temporal window, retrieve a destination event in that window,
  and answer a terminal question about the destination event. This construction
  intentionally isolates interval-conditioned multi-hop retrieval from ordinary
  single-hop lookup, while keeping explicit metadata that allows us to analyze the
  effect of bridge scope, offset, terminal type, distractor type, and surface
  temporal phrasing.

  To reduce benchmark artifacts, generation is not tied to a single template.
  Records are sampled from persona/theme contexts, event-family libraries,
  terminal-question types, bridge scopes, offset ranges, and multiple distractor
  modes. Distractors are constructed to be semantically and structurally close to
  the destination event: for example, they may share the same activity, preserve
  the same object or location skeleton when required by the terminal type, occur
  near the target window, or introduce competing source-like temporal anchors.
  The final judge enforces that the question is not answerable from the source
  alone or the destination alone, but remains uniquely answerable when all source,
  destination, distractor, and filler sessions are present.

  Session generation is separated from event generation. After the event-level
  graph passes the uniqueness gate, each event is realized as a multi-turn memory
  session with a sampled temporal surface form. These forms include deterministic
  calendar references, relative references such as ``last week'' or ``three months
  ago'', and less calendar-like user-memory anchors such as life episodes,
  projects, trips, moves, medical events, habits, or emotionally salient periods.
  This is intended to make the benchmark less dependent on explicit date strings:
  the temporal validator checks whether the realized session still supports the
  intended day, week, or month bucket under the sampled bridge scope.

  We report \textsc{TimeMemEval} as a targeted temporal-bridge retrieval stress
  test. It does not subsume broader temporal phenomena such as long causal chains,
  belief revision, recurring habits, partial temporal order, or contradictory
  memories. To make this limitation explicit, benchmark metadata records the
  temporal topology of every instance, and evaluation can be stratified by bridge
  scope, terminal type, distractor mode, surface temporal style, and whether
  temporal normalization is provided to the model. This allows us to distinguish
  performance gains from temporal normalization, retrieval policy, and reasoning
  over the generated memory evidence, and supports held-out splits over event
  families, surface templates, and temporal relation types.

\subsection{Training and Held-out Evaluation Dataset Construction}
\label{app:train-valid-data-construction}

We construct the GRPO data from two public long-term memory benchmarks,
LoCoMo and LongMemEval, together with our generated \bench records. LoCoMo and
LongMemEval provide real multi-session dialogue questions with gold answers and
evidence-session annotations, while \bench contributes targeted
temporal-causal questions whose answer evidence is not directly addressable from
the original query. These sources play different roles: the public benchmarks
preserve general long-memory coverage, whereas \bench increases the density of
training and evaluation cases that require temporal bridge construction.

\textbf{Source filtering.}
We start from 1,986 LoCoMo questions, 500 LongMemEval questions, and 500 raw
\bench candidates. For LoCoMo, we exclude category-5 adversarial questions; for
LongMemEval, we exclude single-session preference questions, yielding 1,536 and
469 usable public questions, respectively. For \bench, the generation pipeline
first produces 500 raw candidates and then applies automatic validity checks for
invalid generation attempts, duplicate examples, split leakage, incomplete
metadata, temporal inconsistency among target windows, session times and event
times, and answer non-uniqueness. This leaves 374 filtered \bench records that
are eligible for training or held-out evaluation. The resulting filtered pool
contains 2,379 examples.

\textbf{Training and held-out evaluation splits.}
The final GRPO training set contains 500 examples. To avoid conversation-level
leakage in LoCoMo, we first partition LoCoMo by dialogue rather than by
question: six conversations form the LoCoMo training pool and the remaining four
conversations are reserved for held-out evaluation. We then sample 401 LoCoMo
questions from the six training conversations. LongMemEval is not used for GRPO
training; all 469 filtered LongMemEval questions are kept for zero-shot
held-out evaluation. For \bench, we sample 99 training examples from the 374
filtered records and use the remaining 275 filtered records for held-out
evaluation. The held-out evaluation sets are not used for checkpoint selection.

\textbf{Example format.}
Each training example stores only lightweight supervision and metadata: the user
query, gold answer, gold evidence sessions, source identifier, user or speaker
metadata, source bucket, difficulty label, and evaluation mode. It does not store
fixed retrieved memories or precomputed rollout traces. During GRPO, the system
uses the sample metadata to dynamically call the corresponding NanoMem search
environment, generate temporal hints, retrieve sessions, synthesize events,
produce verdicts, and score the final answer. This ensures that the policy is
trained on online memory-search behavior rather than a static retrieved context.

\begin{table}[h]
\centering
\footnotesize
\setlength{\tabcolsep}{3pt}
\begin{tabular}{@{}lrrrrp{0.23\linewidth}@{}}
\toprule
\textbf{Data stage} & \textbf{LoCoMo} & \textbf{LongMemEval} &
\textbf{TME} & \textbf{Total} & \textbf{Use} \\
\midrule
Original questions & 1,986 & 500 & 500 & 2,986 & Source pool \\
Filtered questions & 1,536 & 469 & 374 & 2,379 & Candidate pool \\
Train & 401 & 0 & 99 & 500 & GRPO training \\
Held-out evaluation & 655 & 469 & 275 & 1,399 & Reporting only \\
\bottomrule
\end{tabular}
\caption{Distribution of query pools for GRPO training and held-out evaluation.
LoCoMo is split by conversation; LME is zero-shot; TME uses filtered records.}
\label{tab:train-valid-data-distribution}
\end{table}
